\section{成像、采样与系统前评估：先判断信息能否被采到}
\subsection{这部分为什么应当优先复习}
9月21日第二段录音在 00:54:17 明确举例：小测可以要求计算分辨率需求，也可以回答已讲过的概念题。10月12日 00:00:00--00:00:53 又说明，小测约15分钟、1--2道题，重在理解和思考。因此，复习不能只背相机像素数，应能完成\textbf{任务要求 $\rightarrow$ 最小目标尺寸 $\rightarrow$ 物方采样间隔 $\rightarrow$ 像素数与采样速度 $\rightarrow$ 工程限制}的完整推理。

\noindent\textbf{来源：}9月21日第二段，00:54:17；10月12日，00:00:00、00:02:20。本讲义标注的“课堂案例”指老师确实讨论过的情境；未在录音中给出的数字、推导或条件均另标为补充。

\subsection{图像是什么，系统又包含什么}
单通道静态图像可写成二维函数 $f(x,y)$，其中 $x,y$ 表示位置，函数值表示亮度或测得的某种物理响应。数字化以后，用有限数组 $f[m,n]$ 表示。RGB图像在每个位置有三个分量；视频还增加时间索引 $t$。应区分三件事：\textbf{图像的尺寸}是行列数，\textbf{灰度位深}决定可编码的强度级数，\textbf{物方空间分辨能力}讨论场景中多小的细节可以被区分。三者不是同一个量。

完整系统至少涉及：光源和被测物体、镜头与成像几何、传感器、采样和量化、计算设备、传输、存储、显示或输出。处理链中任何一环不合适，都可能使后面的算法失去意义。例如曝光饱和造成的平坦白区、目标只占一个像素、严重运动模糊，均不能靠调对比度可靠地恢复真实细节。

\textbf{面向应用的三个层次：}发现“这里有物体”、区分“它是什么类别”、精确识别或测量其细小结构，对目标所占像素数的需求不同。能看见车辆不等于能读车牌；能看见裂缝不等于能测准裂缝宽度。像素足够是必要条件之一，通常并非充分条件，还受光学模糊、信噪比、对比度、标定误差影响。

\noindent\textbf{来源：}9月14日，01:54:39及后续道路中线案例；9月21日第二段，00:37:22--00:42附近。

\subsection{FOV与分辨率需求：最重要的一组公式}
设一个方向的视场宽度为 $W$，该方向有效像素数为 $N_x$，物方采样间隔为 $\Delta_x$，则在视场内比例近似均匀时
\begin{equation}
\Delta_x=\frac{W}{N_x},\qquad p=\frac{d}{\Delta_x}=\frac{dN_x}{W}.
\end{equation}
这里 $d$ 是最小目标或缺陷的宽度，$p$ 是其横跨的像素数。若设计要求至少 $k$ 个像素覆盖 $d$，则
\begin{equation}
\boxed{\Delta_x\leq\frac{d}{k},\qquad N_x\geq\left\lceil\frac{kW}{d}\right\rceil .}
\end{equation}
竖直方向同理：$N_y\geq\lceil k_yH/d_y\rceil$。若两方向都受约束，总像素数才是 $N_xN_y$。不能把横向需要的22000像素误写成整幅只需22000像素，也不能只拿总像素数判断横向是否够用。

\textbf{使用条件：}先统一长度单位；明确 $W$ 是实际成像视场而非支架宽度、传送带名义宽度；若有透视，需检查最不利位置的物方分辨率；有效覆盖应扣除裁边并考虑拼接重叠。向上取整只解决整数限制，不代表已经预留工程裕量。

\textbf{补充几何关系：}小视角、近似薄透镜且物距远大于焦距时，若传感器宽度为 $S$、工作距离为 $Z$、焦距为 $f_{\ell}$，有 $W\approx ZS/f_{\ell}$。此关系用于估选镜头；最终应依据真实镜头视场、畸变和安装距离核实。老师强调的思路是先确定需要的场景和距离，再选能满足要求的镜头，而不是先买设备再迁就任务。

\subsection{例1：印刷缺陷检测需要多少横向像素？【课堂案例＋完整计算】}
\textbf{题设与来源：}10月12日 00:02:20，老师讨论包装印刷检测，先说场景宽2.2 m，随后口头简说“两米”；最小缺陷0.5 mm，希望覆盖5个像素。本解明确采用\textbf{实际视场 $W=2.2$ m}，不把两种口述混成一个精确条件。

\textbf{步骤1：统一单位。} $W=2200$ mm，$d=0.5$ mm，$k=5$。

\textbf{步骤2：求物方采样间隔。}
\[
\Delta_x\leq\frac{0.5}{5}=0.1\ \mathrm{mm/pixel}.
\]
\textbf{步骤3：求有效横向像素数。}
\[
N_x\geq\frac{2200}{0.1}=22000.
\]
所以至少需要22000个\textbf{有效横向采样点}。若题目明确改为2.0 m，则答案改为20000，而非22000。

\textbf{步骤4：检验一个候选方案。}补充假设某单相机横向有效像素为16384，则
\[
\Delta_x=\frac{2200}{16384}\approx0.1343\ \mathrm{mm/pixel},\qquad
p=\frac{0.5\times16384}{2200}\approx3.72.
\]
故这个方案不满足“至少5像素”的要求。若用两台相机拼接，不能只相加标称像素数，必须按各自视场计算分辨率，再验证共同覆盖和重叠区域。这里的16384仅用于补充检验，\textbf{不是老师在10月12日这段给出的相机参数}。

\textbf{结论应怎样写：}在2.2 m有效视场、0.5 mm最小缺陷、5像素覆盖的假设下，需要物方采样间隔不超过0.1 mm，横向有效像素数不少于22000；此外还须核实镜头分辨能力、曝光、照明和运动采样。

\subsection{例2：道路中线是否有足够像素？【课堂案例＋一致性检查】}
\textbf{课堂数据：}道路标线宽约15 cm，老师给出的物方像素尺度约4.52 cm/pixel，并强调标线要由多个像素表达，以维持连续性。来源：9月14日道路航拍段落，以及9月21日第二段 00:38:36。
\[
p=\frac{15}{4.52}\approx3.32\ \mathrm{pixels}.
\]
因此标线宽度约覆盖3.3个像素。从“至少3像素”的设计尺度看基本满足；若要求至少5像素，则需要
\[
\Delta_x\leq15/5=3\ \mathrm{cm/pixel},
\]
现有4.52 cm/pixel便不够。

\textbf{重要核验：}录音同时出现7952$\times$5304像素、约300$\times$200 m视场等口述数据。若直接把300 m配到7952像素，则
\[
\frac{30000}{7952}\approx3.77\ \mathrm{cm/pixel},
\]
并不等于4.52 cm/pixel。因此这些近似口述不能同时当成严格一致的一组标定值。本题采用老师直接给出的4.52 cm/pixel来判断标线覆盖；实际题目若明确给视场和像素数，应优先依那组条件计算并说明对应方向。

\textbf{理解题：为什么不只要1像素？}细线与采样网格的位置关系、光学模糊和噪声可能使单像素宽结构断裂或漏检。多像素覆盖提供一定冗余，但“3像素一定不会断”不是普遍保证，仍需结合成像质量评估。

\subsection{运动、行频、曝光和数据量}
线阵相机一次采一行，横向由线阵像素提供采样，纵向由相机与物体的相对运动提供采样。设速度为 $v$，行频为 $f_{\mathrm{line}}$，相邻行的物方间隔为
\begin{equation}
\Delta_y=\frac{v}{f_{\mathrm{line}}},\qquad
f_{\mathrm{line}}\geq\frac{v_{\max}}{\Delta_{y,\max}}.
\end{equation}
行频单位是行/秒，行周期为 $1/f_{\mathrm{line}}$。曝光时间 $t_e$ 期间的运动模糊尺度约为 $b=vt_e$；若允许不超过一个纵向像素，则 $t_e\leq\Delta_y/v$。\textbf{采样频率和曝光时间是两个约束}，行频达标不代表拖影一定合格。

速度变化时，固定行频会使相邻行代表的实际距离变化。编码器可提供位移信息，用来触发行采集或做几何重采样；若纸板相对传送带滑动，单测传送带速度仍不一定等于纸板真实速度。还要检查打滑、机械振动、同步、照明和存储带宽。

\noindent\textbf{来源：}9月28日，00:27:38--00:32:37；10月12日，00:02:20--00:03:07。

\subsection{例3：线阵行频、曝光与带宽【补充练习】}
\textbf{题目：}沿用0.1 mm的目标采样间隔；另设物体最高速度为2 m/s，横向22000像素，单通道8 bit，数据不压缩。求最低行频、允许一个像素拖影时的最长曝光及原始数据率。

\textbf{解：}速度统一为2000 mm/s，因此
\[
f_{\mathrm{line}}\geq2000/0.1=20000\ \mathrm{lines/s},\qquad
T_{\mathrm{line}}\leq50\ \mu\mathrm{s}.
\]
在最高速度下，要使拖影不超过0.1 mm，
\[
t_e\leq0.1/2000=5\times10^{-5}\ \mathrm{s}=50\ \mu\mathrm{s}.
\]
若要求半像素拖影，则最长曝光进一步降至25 $\mu$s。此时往往要增加照明，不能单靠软件把暗噪声变少。

每秒的采样值数为 $22000\times20000=4.4\times10^8$；每个8 bit即1 byte，所以
\[
R=440\ \mathrm{MB/s}=3.52\ \mathrm{Gbit/s}
\]
（采用十进制MB，不含协议、元数据和冗余）。一分钟约26.4 GB。若采集RGB三个8 bit通道，则数据量约为三倍。\textbf{答题重点：}像素需求、行频、曝光、传输和存储必须同时满足，不能只给“买一台高清相机”。

\subsection{采样与量化：不能混淆的两步}
\textbf{采样}使位置（视频还包括时间）离散；\textbf{量化}使幅值离散。$b$ bit可编码 $L=2^b$ 个级别，8 bit灰度通常取0至255，三个8 bit颜色通道对应每像素24 bit。实际数组常以左上角为原点，行号向下增加，列号向右增加；必须说明约定，不能把行列下标和笛卡尔坐标不加区分地互换。

若一维连续信号带限至最高频率 $f_{\max}$，理想无混叠重建要求采样频率高于 $2f_{\max}$（临界相等时还需额外条件）。空间频率若用周期/mm计量，则采样频率用样本/mm计量。二维情况应分别考虑各方向的频谱支持，或使用二维采样格的条件。\textbf{这是理想带限条件下的理论，不能把“2倍”直接解释成任意缺陷只需2像素就一定可靠检测。}

采样不足会混叠，可能出现锯齿、摩尔纹和虚假的低频结构；降低分辨率前通常需低通抗混叠。量化不足易使平滑渐变出现假轮廓、色带。理想采样可在条件成立时恢复连续信号；有限位深量化通常把多个输入幅值合为同一编码，不能逐点精确逆转。上采样和插值增加数组元素数量，不自动增加已采集的真实细节。

\noindent\textbf{来源：}9月14日，语音采样与量化段落；9月21日第二段，00:39:23及灰度级与空间分辨率比较段落。

\subsection{例4：量化误差与存储量【补充练习】}
\textbf{题目：}对归一化强度 $r\in[0,1]$ 进行8 bit端点包含式均匀量化，四舍五入到最近编码。若 $r=0.5$，求量化结果与误差。再计算1920$\times$1080灰度帧和RGB帧的未压缩大小。

\textbf{解：}本题指定重建级别 $\hat r=q/255$，$q\in\{0,\ldots,255\}$，相邻级间隔 $\Delta=1/255$；采用半整数向上取整：
\[
q=\operatorname{round}(255\times0.5)=128,
\quad\hat r=128/255\approx0.5019608,
\quad e=\hat r-r\approx0.0019608.
\]
未过载、取最近级时有 $|e|\leq\Delta/2=1/510$。若题目采用256个等宽量化区间的另一种定义，步长可写为 $1/256$；\textbf{必须先说明量化器约定，不可在同一计算中混用两者}。

像素总数为 $1920\times1080=2073600$。8 bit灰度帧占2073600 byte，约2.074 MB；RGB帧占6220800 byte，约6.221 MB。不含文件头、行填充和压缩。若保持30帧/秒，RGB原始数据率约186.624 MB/s。\textbf{常见错误：}把24 bit写成24 byte，或在没有压缩比条件时把原始大小当作JPEG大小。

\subsection{云端与边缘：可写成简答题的判断框架}
边缘计算在摄像头附近处理，再上传结果；云端或中心端处理可上传图像后集中计算。边缘端通常减小网络传输和远端响应依赖，但需要端侧算力、功耗预算及维护；中心端便于共享算力和统一升级，但受带宽、延迟与网络可靠性限制。哪一种合适由任务约束决定，而非“边缘永远先进”。

\textbf{课堂例：}停车场本地网络较容易部署，可以上传照片到后台识别；广泛分布、移动网络成本高的路侧设备，可在本地识别后上传车牌等结果。正式方案还须考虑隐私、留证、误识别复查以及是否保留原始图像。\textbf{来源：}9月21日第二段，00:32:32--00:36:42。

\section{增强与复原：先说明你究竟要优化什么}
\subsection{两种思路及其边界}
增强关注特定任务下的可见性、对比度或主观观感，可写为 $g=\mathcal E(f)$。例如拉伸灰度、让暗部更明显、突出边缘。复原则从退化模型出发，估计未退化图像，例如
\begin{equation}
f=h*f_0+n,
\end{equation}
其中 $f_0$ 为理想图像，$h$ 为模糊或系统点扩散函数，$n$ 为噪声，$*$ 为卷积。复原利用模型估计 $f_0$；两类方法可能重合，分类重点在\textbf{目标和依据}，不在输出看起来是否更锐。

“知道退化模型就直接求逆”有条件：逆可能不存在、不唯一或对噪声极敏感。以频域形式为例，$F=HF_0+N$，直接相除得到 $F/H=F_0+N/H$；$H$很小的地方噪声会被严重放大，$H=0$处无法如此求逆。因此模型、先验、正则化和噪声水平都重要。生成的合理纹理不等于被验证的原始事实。

\noindent\textbf{来源：}9月21日第一段，00:42:40--00:43:48；9月28日，00:47:26--00:51:27。上式的逆问题条件是为理解课堂“反向处理”而补充的严谨说明。

\subsection{三类常见操作}
\begin{itemize}
\item 点运算：$g(x,y)=T(f(x,y))$，输出只依赖同位置输入强度；坐标不变
\item 邻域运算：输出依赖某一邻域像素，例如平滑、锐化、中值；中值不是线性卷积
\item 变换域运算：先转换表示，在频域等域中处理，再逆变换；线性平移不变滤波可用卷积定理联系空间域和频域
\end{itemize}
评价时应说明用途。适合人观察的增强不必然有利于后续识别，锐度增加不等于信噪比提高，直方图更均匀不等于信息更多。某个像素值变大只能说明它更亮，不能单独证明对比度更好。

\section{灰度变换：会画趋势，更要会解释作用和代价}
\subsection{统一记号和基本变换}
以下 $r$ 表示输入灰度，$s=T(r)$ 表示输出。为避免幂函数尺度混乱，讨论gamma时先使用 $r,s\in[0,1]$。输出到8 bit时再计算 $q=\operatorname{round}(255s)$，并处理越界。
\begin{align}
\text{恒等：}\quad&s=r;\\
\text{反相：}\quad&s=1-r;\\
\text{幂律：}\quad&s=r^\gamma,\qquad\gamma>0;\\
\text{归一化对数：}\quad&s=\frac{\ln(1+ar)}{\ln(1+a)},\qquad a>0.
\end{align}
恒等曲线斜率为正1；反相曲线斜率为负1。不能只凭“45度直线”判断反相，必须看斜率方向。反相将亮暗颠倒；gamma在 $0<\gamma<1$ 时提高中间灰度，$\gamma>1$ 时降低中间灰度，端点0和1保持不变。对数函数在低灰度区斜率较大，可扩展暗部差异并压缩高幅值动态范围，常用于显示幅度跨度大的频谱。

\textbf{局部对比度的解释：}对于相近输入 $r$ 与 $r+\delta$，有
\[
T(r+\delta)-T(r)\approx T'(r)\delta.
\]
$T'(r)>1$ 的区间可扩大原本很小的灰度差；同时也可能放大该区间内的噪声。映射不会区分“有用差异”与“噪声差异”。

\noindent\textbf{来源：}9月28日，01:06:45--01:08:34、01:14:00--01:19:25。这里使用标准公式澄清口头描述中的曲线方向。

\subsection{例5：gamma增强与显示校正【补充练习，源于课堂原理】}
\textbf{题目A：}归一化输入 $r=0.25$，分别施加 $\gamma=0.5,1,2$。
\[
0.25^{0.5}=0.5,\qquad0.25^1=0.25,\qquad0.25^2=0.0625.
\]
转为8 bit并四舍五入，分别为128、64、16。这说明本讲义约定下，gamma小于1提亮，大于1压暗。\textbf{若软件定义的是 $r^{1/\gamma}$，结论会随参数命名反过来；答题先写公式，再谈参数。}

\textbf{题目B：}补充设某显示链路满足 $y=x^{2.2}$，希望最终亮度等于目标 $r$。应先送入什么 $x$？
\[
x=r^{1/2.2},\qquad y=(r^{1/2.2})^{2.2}=r.
\]
若目标 $r=0.25$，则 $x\approx0.5325$；而未经校正直接送0.25，显示亮度约 $0.25^{2.2}\approx0.0474$。校正的思想是与设备响应相补偿。真实显示器色彩管理还包含白点、色域、通道响应与目标传递曲线，不能把课堂简化模型理解为所有设备都只需固定2.2次幂。

\textbf{题目C：}设频谱幅度为0、9、99、999，采用 $s=\ln(1+r)/\ln1000$，输出分别为0、$1/3$、$2/3$、1。原来跨度巨大的幅度被压缩到显示范围，小幅值结构更易观察；这只是显示变换，不说明频谱能量真的被重新分配。

\subsection{线性拉伸、阈值、切片与位平面}
若主要灰度在 $[a,b]$，可采用带截断的拉伸
\begin{equation}
T(r)=\begin{cases}
0,&r\leq a,\\
(L-1)\dfrac{r-a}{b-a},&a<r<b,\\
L-1,&r\geq b.
\end{cases}
\end{equation}
它扩大目标区间，区间外则可能丢失差异。$a=b$时分母为零，必须特别处理，常量图不能靠此法制造新细节。分段线性变换可给特定区间较大斜率、给两端较小斜率，端点位置决定强调哪里。

二值化将输出限制为两种值，例如 $s=0$（$r<\tau$）或 $s=L-1$（$r\geq\tau$）；阈值相等时归入哪类需声明。灰度切片则强调一个区间 $[a,b]$，可令区间内为高值而区间外为0，也可保留区间外原灰度。它与位平面切片不同：前者按强度区间，后者提取二进制的一位。

8 bit灰度可分解为
\[
r=\sum_{j=0}^{7}2^jb_j,\qquad b_j=\left\lfloor\frac{r}{2^j}\right\rfloor\bmod2.
\]
最高位 $b_7=1$ 当且仅当 $r\geq128$。第6位为1的范围是64--127与192--255，不是单个阈值以上全为1。高位通常承载较粗的结构，低位可能含细节和噪声；不能断言所有低位均无信息。

\noindent\textbf{来源：}9月28日，01:19:02--01:25:12；公式与明确的位编号用于消除口述中“第几位”的歧义。

\subsection{例6：三种相近灰度如何显现？【课堂演示情境＋补充数字】}
\textbf{课堂情境：}9月28日01:02:00--01:06:16，老师展示灰度非常接近的隐藏图案，通过拉伸使轮廓显现。录音未给出精确三级数值。\textbf{本题补充设定}三个区域灰度为100、101、102，取 $a=100,b=102,L=256$。
\[
T(100)=0,\qquad T(101)=127.5\longrightarrow128,\qquad T(102)=255.
\]
原来相邻差仅1级，现扩大到约127--128级，因此三个区域变得容易区分。但每个输入灰度仍只对应一个输出灰度，\textbf{没有增加原来三个灰度级承载的信息}。

若灰度100附近有幅度为1的噪声，它也可能被拉伸为很大的亮度跳变。这解释了老师演示中“图案出来，噪声也显现”的现象。若把阈值设为101并规定 $r\geq101$ 为255，则输出为0、255、255，后两个区域合并；二值化适合某些分割目的，却牺牲层次。

附加位平面检查：$r=173=128+32+8+4+1=(10101101)_2$。因此 $b_7=1,b_6=0,b_5=1,b_0=1$；只保留最高两位重建得到128，保留最高四位得到160。这种截断会丢失被删除位所表示的幅值信息。

\section{直方图与均衡化：能计算，也要能指出失效条件}
\subsection{直方图告诉你什么，不能告诉你什么}
对含 $P=MN$ 个像素、灰度级 $r_k=k$ 的图像，定义
\begin{equation}
h(k)=n_k,\qquad p_r(k)=\frac{n_k}{P},\qquad
C(k)=\sum_{j=0}^{k}p_r(j),\qquad \sum_{k=0}^{L-1}p_r(k)=1.
\end{equation}
$h$是计数，$p_r$是归一化频率，$C$是累积分布。横轴是灰度，纵轴是计数或概率，不能混用。直方图集中在低灰度端通常表示整体偏暗；集中在狭窄区间常提示全局对比度不足；双峰可能来自两类亮度区域，但不能仅由双峰证明存在两个具体物体。

直方图丢弃了位置关系。将所有像素随机打乱，直方图完全不变，但形状、纹理和语义可完全改变。全黑背景中一个亮字和同等亮点数的随机散点可以有相同直方图。因此，单靠直方图不能判断清晰度、目标位置、光照空间分布或识别质量；“铺满0--255”也不等于细节丰富。

\noindent\textbf{来源：}9月28日，01:26:21起的直方图讨论，以及01:43:37对“均匀就是最好”的质疑。

\subsection{连续均衡化的推理与离散实现}
对归一化连续灰度随机变量 $R$，假设概率密度为 $p_R(r)$，构造
\begin{equation}
s=T(r)=\int_0^r p_R(w)\,dw.
\end{equation}
在适用的连续、可逆可微区间，概率守恒给出
\[
p_S(s)\,ds=p_R(r)\,dr,\qquad
\frac{ds}{dr}=p_R(r),\qquad p_S(s)=1.
\]
直观理解：按累计概率分配输出位置，原来概率大的区间得到较宽输出跨度。$T$至少单调不减，因此不会颠倒明暗次序。只有在相应区间密度为正时才严格递增，不能笼统断言累积分布处处严格递增。

离散实现的一种明确约定是
\begin{equation}
\boxed{s_k=\operatorname{round}\big[(L-1)C(k)\big].}
\end{equation}
先统计、再累计、再建立查找表、最后逐像素替换；输出直方图必须把映射到同一输出级的计数相加。有限像素、整数灰度和确定性的逐级映射，使输出通常\textbf{不可能严格均匀}：同一级原像素被整体映到同一输出值，不能任意拆分。不同源级也可能因取整合并。

某些实现用 $C_{\min}$修正：$s_k=\operatorname{round}[(L-1)(C(k)-C_{\min})/(1-C_{\min})]$，其中 $C_{\min}$ 为首个非零累计概率。它与本讲义主公式不是同一个约定；遇到常量图 $C_{\min}=1$ 需特殊处理。\textbf{手算时写清公式与取整，不混用实现版本。}

\noindent\textbf{来源：}9月28日，01:35:27--01:41:59；课堂说明理解累计变换思路即可，不强调重新推导全套微积分。

\subsection{例7：完整离散均衡化手算【补充练习】}
\textbf{题目：}8级灰度图像有16个像素，0至7的计数依次为
\[
(n_0,\ldots,n_7)=(0,2,4,6,2,2,0,0).
\]
采用 $s_k=\operatorname{round}[7C(k)]$，半整数向上取整，计算映射与输出直方图。

\textbf{步骤1：核对总数与概率。}计数和为16，故概率为
\[
(0,1/8,1/4,3/8,1/8,1/8,0,0).
\]
\textbf{步骤2：累加并缩放。}
\[
\begin{array}{c|rrrrrrrr}
k&0&1&2&3&4&5&6&7\\\hline
n_k&0&2&4&6&2&2&0&0\\
16C(k)&0&2&6&12&14&16&16&16\\
7C(k)&0&0.875&2.625&5.25&6.125&7&7&7\\
s_k&0&1&3&5&6&7&7&7
\end{array}
\]
\textbf{步骤3：把计数搬到输出级。}有效映射是 $1\mapsto1$、$2\mapsto3$、$3\mapsto5$、$4\mapsto6$、$5\mapsto7$。于是
\[
(n'_0,\ldots,n'_7)=(0,2,0,4,0,6,2,2).
\]
输出仍为16个像素。原来占用1--5，输出占用1--7，灰度跨度扩展，但直方图\textbf{并不均匀}，原来灰度3处的6个像素仍构成一个6像素的大柱。

\textbf{步骤4：回答概念追问。}本例非空灰度数仍为5，说明均衡化不一定减少有效灰度级；课堂强调的减少是离散映射发生合并时可能出现的代价。全局点映射不会增加原本可区分的灰度类别；若两个源级被合并，就不能凭输出唯一恢复它们。

\textbf{与课堂例子的关系：}9月28日01:40:31讲过一个64$\times$64的例子，口述映射为 $0\mapsto1$、$1\mapsto3$、$2\mapsto5$、$3,4\mapsto6$、$5,6,7\mapsto7$。该例说明8个源级可能只剩5个输出级；本题使用独立构造计数，不把不完整口述补成原题。

\subsection{直方图规定化与局部增强}
若想让输出接近指定分布，而非均匀分布，设源CDF为 $T(r)$、目标CDF为 $G(z)$，连续情况下使用
\begin{equation}
z=G^{-1}(T(r)).
\end{equation}
逻辑是源灰度先映到累计概率，再在目标CDF中寻找对应灰度。离散CDF可能没有严格反函数，需指定“最接近累计概率”或“首个不小于该概率”的规则；不同规则可能得出不同近似结果。确定性灰度映射通常不能保证精确达到任意目标直方图。

全局均衡只使用全图统计。若亮暗不均分布在不同区域，大面积背景可能主导映射，局部的小目标仍不可见。局部直方图均衡在窗口或分块内建立映射；窗口小则更强调局部细节，但也更容易放大噪声，产生不稳定映射。分块独立处理还可能出现边界接缝，需平滑衔接。

CLAHE的基本思路是限制局部直方图柱的高度、重新分配被裁剪的计数，再结合相邻块的映射以缓和块边界。它能控制过度增强，但不保证所有场景最佳；块大小和裁剪限制仍要依据对象尺度、噪声和任务选择。\textbf{来源：}9月28日，01:44:30--01:46:28、01:54:29附近；10月12日结尾对局部不均匀与先增强后去噪的讨论。

\subsection{例8：直方图匹配为什么只能近似？【补充练习】}
\textbf{题目：}四级源图像的概率为 $(1/2,1/4,1/4,0)$，希望目标概率为 $(1/4,1/4,1/4,1/4)$。采用“选择使目标CDF与源CDF差值最小的输出级”规则，求匹配。

\textbf{解：}源CDF是 $(1/2,3/4,1,1)$，目标CDF是 $(1/4,1/2,3/4,1)$。对应查找得到
\[
0\mapsto1,\qquad1\mapsto2,\qquad2\mapsto3,
\]
源级3没有像素，映射到哪里不影响结果。因此输出概率为 $(0,1/2,1/4,1/4)$，并未达到所要求的均匀分布。

\textbf{为什么？}源级0含有一半像素，单值灰度变换必须把它们统一映到同一级，不能一半放0、另一半放1。若允许按位置或排序拆分相同灰度像素，则可以实现另外一类更精确的直方图分配，但这已经改变了“同灰度同输出”的点映射规则。故不能把连续公式的理论效果直接套到有限离散数据。

\textbf{延伸简答：一幅整体不错、只有暗角细节不足的图，是否立刻做全局均衡？}不应无条件采用。先检查暗角是否存在有效信号及噪声，比较局部增强或针对照明的处理；若已饱和或没有足够采样信息，直方图方法无法补回真实细节。根据目标区域可见性、噪声和伪影评价，而非只看输出直方图是否更平。

\subsection{这一部分的高频失分点与答题模板}
\begin{enumerate}
\item 把“有更多像素”说成“真实分辨率一定更高”：须同时考虑视场、光学与采样信息
\item 用 $2f_{\max}$机械代替工业任务的多像素覆盖要求：两者约束对象不同
\item 说gamma大于1必定提亮：先写采用 $r^\gamma$ 还是 $r^{1/\gamma}$
\item 把反相说成正斜率45度线：反相应为负斜率；恒等才是正斜率
\item 说均衡化一定均匀、一定更好、一定减少级数：三个“一定”都不成立
\item 只统计输出映射不合并计数：均衡化最终直方图需按输出级重新累计
\item 把灰度切片当成位平面：区间强调与提取二进制位是不同运算
\item 忽略噪声与不可逆丢失：增强会同时改变细节和噪声，截断与级数合并不能简单逆转
\end{enumerate}

\textbf{15分钟小测的简答组织：}先写目标与条件，再写公式或模型，代入统一单位的数值，最后用一两句话说明结论的适用范围。若题目问“为什么某算法不好”，按“图像特点 $\rightarrow$ 算法依据 $\rightarrow$ 不匹配之处 $\rightarrow$ 更合适的方向”回答，避免只列算法名称。

